\documentclass[11pt,a4paper]{article}
\usepackage[T1]{fontenc}
\usepackage{lmodern}
\usepackage[top=0.38in, bottom=0.38in, left=0.68in, right=0.68in, headheight=0pt, headsep=0pt]{geometry}
\usepackage{amsmath,amssymb}
\usepackage{xcolor}
\usepackage{tikz}
\usetikzlibrary{arrows.meta,positioning,calc,decorations.pathreplacing}
\usepackage{tcolorbox}
\usepackage{hyperref}
\usepackage{booktabs}
\usepackage{caption}

\hypersetup{
    colorlinks=true,
    linkcolor=blue!70!black,
    citecolor=blue!70!black,
    urlcolor=blue!70!black
}

% Custom Callout Boxes matching house style
\newtcolorbox{learningobj}{
    colback=blue!3!white,
    colframe=blue!60!black,
    fonttitle=\bfseries\sffamily,
    fontupper=\small,
    title={$\blacktriangleright$ LEARNING OBJECTIVE},
    arc=1.2mm,
    boxrule=0.7pt,
    left=2.5mm, right=2.5mm, top=1.0mm, bottom=1.0mm,
    before skip=1.0mm, after skip=1.5mm
}

\newtcolorbox{groundbreaking}{
    colback=teal!4!white,
    colframe=teal!60!black,
    fonttitle=\bfseries\sffamily,
    fontupper=\small,
    title={$\bigstar$ GROUNDBREAKING FINDING},
    arc=1.2mm,
    boxrule=0.7pt,
    left=2.5mm, right=2.5mm, top=1.0mm, bottom=1.0mm,
    before skip=1.0mm, after skip=1.5mm
}

\newtcolorbox{physicalinsight}{
    colback=yellow!6!white,
    colframe=orange!75!black,
    fonttitle=\bfseries\sffamily,
    fontupper=\small,
    title={$\blacktriangleright$ PHYSICAL INSIGHT},
    arc=1.2mm,
    boxrule=0.7pt,
    left=2.5mm, right=2.5mm, top=1.0mm, bottom=1.0mm,
    before skip=1.0mm, after skip=1.5mm
}

\setlength{\parindent}{0pt}
\setlength{\parskip}{2.0pt}

\newcommand{\diagsec}[1]{\vspace{3pt}\textbf{\large #1}\vspace{1pt}\par}

\begin{document}
\setlength{\abovedisplayskip}{1.5pt}
\setlength{\belowdisplayskip}{1.5pt}

% ==================== PAGE 1 ====================
\subsection*{1.6 From Recognition to Evaluation: Two Algorithms, Not One}

\begin{learningobj}
Rather than treating speech evaluation as a downstream byproduct of standard Automatic Speech Recognition (ASR), this section deconstructs the fundamental algorithmic divide between statistical transcription and acoustic verification. By dissecting the failure modes of linear pitch-tracking baselines on physical speech, we derive the mathematical and biomechanical necessity of dual-tier Dynamic Time Warping (DTW) for real-time edge evaluation.
\end{learningobj}

\textbf{Organic Bridge from Section 1.5: The Boundary Between Recognition and Verification.} In Section~1.5, we stress-tested the front-end DSP pipeline against physical silicon and acoustic boundaries, proving that nominal parameters ($16\text{ kHz}$ sampling, $25\text{ ms}$ windowing, $N=512$ zero-padding, and sparse INT16 Mel compression) represent an immovable Pareto-optimal operating point. That front-end emits an 80-dimensional Log-Mel energy vector every $10\text{ ms}$.

In conventional voice architectures, these spectral feature vectors feed directly into an Automatic Speech Recognition (ASR) acoustic model (such as a Conformer or CTC decoder) to transcribe spoken audio into written text. This pipeline operates under a fundamental probabilistic objective:
\begin{equation}
  \hat{W} = \arg\max_W P(W \mid \mathbf{X}) \propto P(\mathbf{X} \mid W) \cdot P(W)
\end{equation}
where $P(W)$ is a statistical language model prior. In speech recognition, the language model prior is explicitly engineered to be forgiving: if a non-native speaker mispronounces a vowel or slurs a consonant, the language model calculates that the erroneous acoustic token has near-zero probability in human language, silently autocorrecting the transcription into the most likely grammatical word. In ASR, forgiving human speech errors is a primary design virtue.

In educational voice AI, pronunciation training, and clinical speech assessment, however, this exact virtue becomes a catastrophic failure. The goal of an evaluation engine is not to guess what the user meant to say; the target reference text is known 100\% in advance. The task is to measure precisely \emph{how} the user articulated the speech sounds compared to an authoritative reference. If an evaluation system utilizes standard ASR decoders, the language model prior hallucinates correctness, awarding high confidence scores to misarticulated phonemes. 

Speech evaluation requires a fundamentally different algorithmic class: \textbf{deterministic acoustic trajectory verification}. Rather than maximizing posterior token probabilities, the system must act as an unforgiving physical mirror, comparing the user's acoustic spectrogram and pitch melody against a reference standard.

\diagsec{Box 1 --- The Group Baseline: What 492026 Actually Built}

To ground this engineering challenge in concrete hardware reality, we examine the baseline pronunciation evaluation system engineered by the student research team (documented in project report \emph{B{\'A}O C{\'A}O 492026.pdf}, dated 04--23/09/2026 by \DJ{\'u}c Nguy{\~{\^e}}n Ti{\'{e}}n). The baseline system attempted to evaluate English pronunciation prosody on embedded hardware by executing a four-stage algorithmic pipeline:
\begin{enumerate}\setlength{\itemsep}{0pt}
  \item \textbf{Pitch Contour Extraction:} Extract fundamental frequency ($F_0$) trajectories of both reference teacher audio and student learner audio using autocorrelation-based pitch tracking (Praat algorithm).
  \item \textbf{Linear Temporal Normalization:} Because the student speaks at a different tempo than the teacher ($N_{\text{stu}} \ne N_{\text{ref}}$), uniformly stretch or compress the student's pitch vector to match the teacher's length using linear interpolation (\texttt{np.interp}).
  \item \textbf{Sliding Window Segmentation:} Divide the resampled pitch trajectory into $K+1$ overlapping inspection windows to assess intonation across individual words.
  \item \textbf{Prosodic Correlation Scoring:} Compute Pearson's correlation coefficient ($r$) between the normalized pitch contours within each window.
\end{enumerate}

\clearpage
% ==================== PAGE 2 ====================
\textbf{Mathematical Rectification: The Period vs. Analysis Window Mistake.} In documenting the pitch extraction engine, the original report stated: $\text{“Pitch floor } f_{\min} = 75\text{ Hz} \implies T = 1/75 = 0.04\text{ s (40 ms).”}$ This equation conflates two fundamentally different physical quantities: the \textbf{fundamental pitch period} ($T_0$) and the \textbf{autocorrelation analysis window duration} ($T_{\text{win}}$). The fundamental period of a $75\text{ Hz}$ acoustic wave is strictly:
\begin{equation}
  T_0 = \frac{1}{f_{\min}} = \frac{1}{75\text{ Hz}} = 0.013333\text{ s} \quad (13.33\text{ ms})
\end{equation}
If an embedded designer were to configure an analysis window of length $25\text{ ms}$ (as used in standard STFT front-ends), the window would span only $25\text{ ms} / 13.33\text{ ms} \approx 1.875$ pitch cycles. An autocorrelation lag engine requires at least one complete baseline cycle, a full time-shift lag, and a second complete cycle to compute a cross-correlation peak. With fewer than two complete cycles, the autocorrelation function cannot reliably resolve the fundamental period, causing octave jumping errors (mistaking $75\text{ Hz}$ for $150\text{ Hz}$). To guarantee mathematical stability, Praat enforces the \textbf{Three-Period Rule}: the analysis window must span at least three complete glottal cycles of the lowest detectable frequency:
\begin{equation}
  T_{\text{win}} = \frac{3}{f_{\min}} = \frac{3}{75\text{ Hz}} = 0.040\text{ s} \quad (40\text{ ms})
\end{equation}
The value $0.04\text{ s}$ is not the wave period; it is the physical window integration floor required to detect a $75\text{ Hz}$ glottal oscillation.

\textbf{Worked Pitch Extraction Count.} Consider an actual reference recording from teacher ``Sarah'' with duration $T = 1.450667\text{ s}$. Under Praat's standard configuration with analysis window $T_{\text{win}} = 0.040000\text{ s}$ ($40\text{ ms}$) and hop cadence $\Delta t = 0.010000\text{ s}$ ($10\text{ ms}$, advancing at $100\text{ Hz}$), the active temporal span available for pitch evaluation without zero-padding is:
\begin{equation}
  T_{\text{active}} = T - T_{\text{win}} = 1.450667\text{ s} - 0.040000\text{ s} = 1.410667\text{ s}
\end{equation}
Dividing by the $10\text{ ms}$ hop step yields $\lfloor 1.410667\text{ s} / 0.010000\text{ s} \rfloor = 141$ intervals, producing exactly $141 + 1 = 142\text{ discrete pitch points}$. The residual fraction $\Delta t_{\text{residual}} = 1.450667\text{ s} - (141 \times 0.01\text{ s} + 0.04\text{ s}) = 0.000667\text{ s}$ ($667\ \mu\text{s}$) is symmetrically split into two boundary margins of $0.000333\text{ s}$ ($333\ \mu\text{s}$) centered at the recording extremities.

\textbf{Uniform Linear Resampling Quantization Failure.} When comparing Sarah's reference audio ($T_1 = 1.42\text{ s}, n_1 = 139\text{ frames}$) against a student attempt ($T_2 = 2.00\text{ s}, n_2 = 196\text{ frames}$), linear interpolation calculates an effective fractional step size:
\begin{equation}
  \Delta t_{\text{resample}} = \frac{T_1}{n_2} = \frac{1.42\text{ s}}{196} \approx 0.00724489\text{ s}
\end{equation}
Fixed-point rounding and frame quantization in embedded software truncate this step to $0.007\text{ s}$. Over 196 frames, the cumulative evaluated duration reaches only $196 \times 0.007\text{ s} = 1.372\text{ s}$. Subtracting this from physical duration reveals an unmodeled boundary gap of $\Delta T_{\text{missed}} = 1.420\text{ s} - 1.372\text{ s} = 0.048\text{ s}$ ($48\text{ ms}$). At a $10\text{ ms}$ frame cadence, this boundary quantization gap discards $\approx 4.8$ ($4\text{--}5$) critical pitch frames at the terminal word boundary, clipping sentence-final intonation inflections.

\textbf{The Heuristic Sliding Window Invariant.} To evaluate prosodic contour locally across words without manual time-alignment, the group derived a mathematical window sizing rule. To construct $K+1$ inspection windows of uniform duration $W$ frames advancing with a $50\%$ overlap ($S = W/2$) over total duration $N$ frames: $N = W + K \times S = W + K \times (W/2) = W(K+2)/2 \implies W = 2N / (K+2)$. The report assumed equal windows; this relation is a mathematical sizing identity following directly from $S = W/2$ and $K+1$ windows, not a measured acoustic or linguistic word boundary.

\textbf{The Group Condemnation Rule.} In the report's scoring protocol, if an analysis window produces $r < 0.3$, every word overlapping that window is condemned and marked as mispronounced---even if the same word also sits on an adjacent overlapping window with $r \ge 0.3$. Because windows overlap by $50\%$, a single local phase misalignment, vowel hesitation, or unvoiced gap within one window cascades to condemn correctly spoken words across the sentence.

\begin{physicalinsight}
The sizing identity $W = 2N / (K + 2)$ treats speech as a sequence of identical, rigid geometric blocks. In human language, word durations vary wildly based on lexical stress, phonetic makeup, and conversational emphasis: the article ``a'' may last $80\text{ ms}$, while the word ``extraordinary'' lasts $750\text{ ms}$. Forcing a uniform window width $W$ slices words arbitrarily in half. Crucially, when an unvoiced stop consonant (/t/, /p/, /k/) occurs, pitch is physically non-existent ($F_0 = 0$). In a rigid sliding window, the zero-pitch void of an unvoiced consonant in word A spills into word B, destroying the variance calculation and triggering cascading false failures.
\end{physicalinsight}

\clearpage
% ==================== PAGE 3 ====================
\diagsec{Box 2 --- Where Linear Pitch Pearson Breaks: The Biomechanical Collision}

The fatal flaw of the group's baseline system lies not in the pitch extraction parameters, but in the core assumption of \textbf{Uniform Linear Resampling}.

When a beginner student speaks slowly or hesitates, how does the physical sound stretch?
\begin{itemize}\setlength{\itemsep}{0pt}
  \item \textbf{Hypothesis of the Naive Engineer:} If a student takes $1.8\text{ s}$ to speak a phrase that the teacher spoke in $1.0\text{ s}$ (an $80\%$ increase in duration), every millisecond, phoneme, vowel, and consonant was stretched by exactly $1.8\times$. Time can be normalized using an \textbf{iron ruler}.
  \item \textbf{Physical and Biomechanical Reality:} Human speech is an \textbf{elastic rubber band}, not an iron ruler. 
\end{itemize}

\textbf{Biomechanical Asymmetry: Vowels vs. Stop Consonants.}
\begin{enumerate}\setlength{\itemsep}{0pt}
  \item \textbf{Vowels are aerodynamic resonances:} During a vowel sound (/a/, /i/, /u/), the vocal tract is open and unobstructed. The lungs continue to supply air, and the vocal cords vibrate in steady-state resonance. A speaker can sustain a vowel for $200\text{ ms}$, $1\text{ second}$, or $5\text{ seconds}$ at will.
  \item \textbf{Stop consonants are ballistic releases:} Producing a stop consonant (/d/, /t/, /b/, /p/) requires building intra-oral air pressure behind a complete closure (tongue or lips) and releasing it. The physical duration of the acoustic burst is determined strictly by the biomechanical inertia of the lips and tongue muscles: typically $20\text{ ms}$ to $35\text{ ms}$. A speaker physically \textbf{cannot} stretch a /t/ burst to $200\text{ ms}$---attempting to do so produces silence followed by a normal burst.
\end{enumerate}
When a student hesitates, they stretch vowels elastically while leaving stop consonants unextended. Uniform linear resampling compresses the prolonged vowel across subsequent consonants and words, shifting acoustic boundaries out of phase.

\textbf{The Rise/Fall Phase Inversion Disaster.} Consider a simple two-word pedagogical assessment:
\begin{itemize}\setlength{\itemsep}{0pt}
  \item \textbf{Teacher Reference:} ``Good morning'' ($1.0\text{ s}$, 100 frames).  
  ``Good'' is articulated in $0.35\text{ s}$ ($35\%$ of duration) with a rising pitch contour ($160 \to 190\text{ Hz}$). ``morning'' is articulated in $0.65\text{ s}$ ($65\%$ of duration) with a falling pitch cadence ($175 \to 120\text{ Hz}$).
  \item \textbf{Student Utterance:} ``Gooood... morning'' ($1.8\text{ s}$, 180 frames).  
  The student hesitates on ``Good'', sustaining the /u/ vowel for $1.1\text{ s}$ ($61\%$ of duration) with a rising pitch ($160 \to 190\text{ Hz}$). ``morning'' is spoken normally in $0.7\text{ s}$ ($39\%$ of duration) with a falling pitch ($175 \to 120\text{ Hz}$).
\end{itemize}
Both speakers executed the exact same melodic contour: rise on ``Good'', fall on ``morning''.

Now observe what happens when uniform linear resampling compresses the student's 180 frames into 100 frames:
\begin{itemize}\setlength{\itemsep}{0pt}
  \item In the teacher reference, frame 45 is $45\%$ into the sentence, which lies inside ``morning''. The teacher's pitch is falling: $\Delta y_{45} = y_{45} - \bar{y} < 0$.
  \item In the student's resampled vector, because ``Gooood...'' occupied $61\%$ of the original utterance, frame 45 is \textbf{still inside the prolonged word ``Good''}! The student's pitch is rising: $\Delta x_{45} = x_{45} - \bar{x} > 0$.
\end{itemize}
When the system calculates the Pearson correlation covariance term at frame 45:
\begin{equation}
  \text{Covariance Term} = \Delta x_{45} \cdot \Delta y_{45} = (+\text{rising}) \times (-\text{falling}) < 0
\end{equation}
Instead of comparing ``Good'' against ``Good'' and ``morning'' against ``morning'', the iron ruler forces the rising melody of ``Good'' to multiply the falling melody of ``morning''. Across the entire evaluation window, negative products dominate the numerator, collapsing the correlation to $r \approx -0.42$. The system rejects the student's attempt, reporting ``Incorrect Intonation'', despite the student having mimicked the melody flawlessly.

\clearpage
% ==================== PAGE 4 ====================
\textbf{Worked 4-Point Mathematical Breakdown: Constructed Contrary-Motion Probe.} To make the algebraic mechanism irrefutable, we examine a constructed contrary-motion probe demonstrating how opposite pitch trajectories collapse into a severe negative correlation:
\begin{itemize}\setlength{\itemsep}{0pt}
  \item \textbf{Teacher Pitch Contour:} $\mathbf{x} = [120, 140, 160, 130]\text{ Hz}$. Mean: $\bar{x} = \frac{120 + 140 + 160 + 130}{4} = 137.5\text{ Hz}$. Centered deviations: $\Delta \mathbf{x} = [-17.5, +2.5, +22.5, -7.5]\text{ Hz}$. Variance sum: $\sum (\Delta x_i)^2 = (-17.5)^2 + (2.5)^2 + (22.5)^2 + (-7.5)^2 = 875.0$.
  \item \textbf{Student Pitch Contour:} $\mathbf{y} = [140, 120, 110, 150]\text{ Hz}$. Mean: $\bar{y} = \frac{140 + 120 + 110 + 150}{4} = 130.0\text{ Hz}$. Centered deviations: $\Delta \mathbf{y} = [+10.0, -10.0, -20.0, +20.0]\text{ Hz}$. Variance sum: $\sum (\Delta y_i)^2 = 10^2 + (-10)^2 + (-20)^2 + 20^2 = 1{,}000.0$.
\end{itemize}

Now compute the cross-product terms $\Delta x_i \cdot \Delta y_i$:

\begin{table}[h!]
\centering\small
\caption*{Table 1.6: Constructed Contrary-Motion Probe}
\begin{tabular}{ccccccl}
\toprule
\textbf{Index $i$} & $x_i$ (Hz) & $\Delta x_i$ & $y_i$ (Hz) & $\Delta y_i$ & $\Delta x_i \cdot \Delta y_i$ & \textbf{Physical Contribution} \\
\midrule
1 & 120 & $-17.5$ & 140 & $+10.0$ & $-175.0$ & Contrary motion (Negative penalty) \\
2 & 140 & $+2.5$  & 120 & $-10.0$ & $-25.0$  & Contrary motion (Negative penalty) \\
3 & 160 & $+22.5$ & 110 & $-20.0$ & $-450.0$ & Peak misaligned with valley (Severe penalty) \\
4 & 130 & $-7.5$  & 150 & $+20.0$ & $-150.0$ & Contrary motion (Negative penalty) \\
\midrule
\textbf{Sum} & & & & & \textbf{$-800.0$} & \textbf{Net Negative Covariance} \\
\bottomrule
\end{tabular}
\end{table}

Evaluating Pearson's correlation coefficient:
\begin{equation}
  r = \frac{\sum \Delta x_i \Delta y_i}{\sqrt{\sum (\Delta x_i)^2} \cdot \sqrt{\sum (\Delta y_i)^2}} = \frac{-800.0}{\sqrt{875.0} \cdot \sqrt{1000.0}} = \frac{-800.0}{29.5804 \times 31.6228} \approx -0.8552
\end{equation}
The result is a severe negative correlation ($r = -0.86$). Opposite pitch motions force a mathematically sound correlation metric to produce the exact opposite of pedagogical truth.

\textbf{Two Fatal Acoustic Blind Spots of Pitch-Only Evaluation.}
\begin{enumerate}\setlength{\itemsep}{0pt}
  \item \textbf{The Humming Trick (Formant Blindness):} The fundamental frequency $F_0$ measures only the rate of vocal fold vibration. It carries zero information regarding oral tract articulation (formants $F_1, F_2, F_3$). If a user keeps their mouth shut and hums (\emph{``Mmm... mmm-mmm''}), their pitch matches the teacher, achieving a near-perfect intonation match without articulating a single phoneme! A closed-mouth hum can match $F_0$ and still miss every formant resonance.
  \item \textbf{The Unvoiced Abyss:} Voiceless consonants (/s/, $/\int/$, /t/, /p/, /k/) involve zero vocal fold vibration ($F_0 = 0$). Linear interpolation connects voiced pitch values across unvoiced regions with artificial slopes, creating fictitious melodic cliffs that distort correlation scores.
\end{enumerate}

\diagsec{Box 3 --- The Option Space: Three Ways to Fix It}

\begin{table}[h!]
\centering
\resizebox{\textwidth}{!}{
\begin{tabular}{lp{4.8cm}p{3.8cm}p{5.4cm}}
\toprule
\textbf{Paradigm} & \textbf{Feature \& Resampling Mechanism} & \textbf{Silicon Profile} & \textbf{Vulnerability / Failure Mode} \\
\midrule
\textbf{Linear pitch Pearson} & Fundamental frequency ($F_0$), uniform global resample (\texttt{np.interp}). & Minimal compute; 1-d interpolation and scalar dot products. & Blind to local tempo variations and unvoiced gaps ($F_0=0$); phase inversion under vowel prolongation. \\
\addlinespace[2pt]
\textbf{Windowed pitch Pearson} (Group baseline) & Pitch ($F_0$) in $K+1$ windows ($S = W/2$), linear resample per window. & Low compute; evaluated independently per window. & Still linear inside each window; delayed fall still multiplies a rise ($\Delta x \cdot \Delta y < 0$); harsh window condemnation ($r < 0.3$). \\
\addlinespace[2pt]
\textbf{Mel-DTW, then warped Pearson} (Our Architecture) & 80-d Log-Mel path $\mathcal{P}$; pitch ($F_0$) reindexed along $\mathcal{P}$. & Balanced 2-d dynamic programming; maps to BRAM and systolic arrays. & Without a Sakoe--Chiba band, or even with one, DTW does not prove the student said the words; gibberish inside band can still align. \\
\bottomrule
\end{tabular}
}
\end{table}

\clearpage
% ==================== PAGE 5 ====================
\diagsec{The Mel-DTW Elastic Rescue}

\textbf{Mathematical Formulation.} Let $\mathbf{X} = [\mathbf{x}_1, \dots, \mathbf{x}_N]$ be the student's 80-dimensional Log-Mel spectrogram sequence ($N$ frames), and $\mathbf{Y} = [\mathbf{y}_1, \dots, \mathbf{y}_M]$ be the teacher's reference sequence ($M$ frames).
\begin{enumerate}\setlength{\itemsep}{0pt}
  \item \textbf{Local Spectral Cost Matrix:} The local acoustic distance between student frame $i$ and reference frame $j$ is the squared Euclidean distance across all 80 Mel frequency bins:
  \begin{equation}
    d(i, j) = \sum_{k=1}^{80} (x_{i,k} - y_{j,k})^2
  \end{equation}
  \item \textbf{Bellman Recurrence Relation:} The optimal cumulative acoustic alignment cost $D(i, j)$ is computed via dynamic programming:
  \begin{equation}
    D(i, j) = d(i, j) + \min \Big\{ D(i-1, j-1), \; D(i-1, j), \; D(i, j-1) \Big\}
  \end{equation}
  The three transitions represent the physical mechanics of speech articulation:
  \begin{itemize}\setlength{\itemsep}{0pt}
    \item \textbf{Diagonal Step $(i-1, j-1)$:} Synchronized phonetic progress at matching tempo.
    \item \textbf{Vertical Step $(i-1, j)$:} Student vowel prolongation or pause (student advances, reference stationary).
    \item \textbf{Horizontal Step $(i, j-1)$:} Student rushing or reference sound deletion (student stationary, reference advances).
  \end{itemize}
  \item \textbf{Global Corridor Constraint:} To prevent pathological alignments, search is bounded within a Sakoe-Chiba band of width $W_{\text{band}}$: $|i - j| \le W_{\text{band}}$.
\end{enumerate}

\textbf{Worked $4 \times 5$ Dynamic Programming Grid.} To observe how dynamic programming elastically absorbs vowel prolongation, consider a concrete numerical grid with $N = 4$ student frames and $M = 5$ reference frames.

Below is the precomputed 80-d local spectral distance matrix $d(i, j)$:
\begin{equation}
\begin{array}{c|ccccc}
d(i, j) & j=1 & j=2 & j=3 & j=4 & j=5 \\ \hline
i=1 & 2 & 4 & 7 & 6 & 3 \\
i=2 & 5 & 3 & 2 & 5 & 4 \\
i=3 & 8 & 6 & 1 & 3 & 6 \\
i=4 & 9 & 7 & 4 & 2 & 1 \\
\end{array}
\end{equation}

Applying the Bellman recurrence with boundary condition $D(1, 1) = d(1, 1) = 2$ and outer edges initialized to $\infty$ yields the cumulative cost matrix $D(i, j)$:
\begin{equation}
\begin{array}{c|ccccc}
D(i, j) & j=1 & j=2 & j=3 & j=4 & j=5 \\ \hline
i=1 & \mathbf{2} & 6 & 13 & 19 & 22 \\
i=2 & 7 & \mathbf{5} & \mathbf{7} & 12 & 16 \\
i=3 & 15 & 11 & 6 & \mathbf{9} & 15 \\
i=4 & 24 & 18 & 10 & 8 & \mathbf{9} \\
\end{array}
\end{equation}

Tracing backward from terminal state $(4, 5)$ along the minimal predecessor cells establishes the optimal warping path $\mathcal{P}$:
\begin{equation}
  (1, 1) \longrightarrow (2, 2) \longrightarrow (2, 3) \longrightarrow (3, 4) \longrightarrow (4, 5)
\end{equation}

\textbf{Critical Acoustic Observation at Frame $i=2$:} At frame $i=2$, the path executes a horizontal transition $(2, 2) \to (2, 3)$: the reference frame advanced while the student frame was held. The dynamic programming grid elastically absorbs this temporal difference without distorting subsequent alignments, ensuring that phonemes at $i=3$ and $i=4$ align precisely with reference phonemes $j=4$ and $j=5$. Boundary alignment is preserved without phase shift.

\clearpage
% ==================== PAGE 6 ====================
\textbf{Rescuing Pearson: DTW-Guided Pitch Evaluation.} With the optimal warping path $\mathcal{P} = \{ (i_k, j_k) \}$ ($k=1, \dots, P$) extracted by Tier~1 spectral DTW, we evaluate pitch prosody along this non-linear path:
\begin{equation}
  \tilde{x}_k = X[i_k], \qquad \tilde{y}_k = Y[j_k] \quad \text{for } k = 1, 2, \dots, P
\end{equation}
Filtering out mutually unvoiced points ($F_0 = 0$), we compute Pearson correlation over the warped pitch vectors:
\begin{equation}
  r_{\text{warped}} = \frac{\sum_{k=1}^P (\tilde{x}_k - \bar{\tilde{x}})(\tilde{y}_k - \bar{\tilde{y}})}{\sqrt{\sum_{k=1}^P (\tilde{x}_k - \bar{\tilde{x}})^2} \cdot \sqrt{\sum_{k=1}^P (\tilde{y}_k - \bar{\tilde{y}})^2}}
\end{equation}
Because identical phonetic moments are compared against each other, phase inversion is completely eliminated. As illustrated in Figure~1.6a, this dual-tier architecture rescues the evaluation score from a false negative failure ($r < 0$) into verified in-phase correlation ($r_{\text{warped}} > 0$).

\textbf{Computational Workload and The Silicon Handoff.} For a typical 5-second spoken phrase evaluated at a $10\text{ ms}$ hop cadence, both student and reference sequences span $N = M = 500\text{ frames}$. Evaluating the unconstrained dynamic programming grid requires:
\begin{equation}
  \text{Total Cells} = 500 \times 500 = 250{,}000\text{ grid points}
\end{equation}
At each cell, the processor must evaluate:
\begin{itemize}\setlength{\itemsep}{0pt}
  \item 80 subtractions and 80 multiply-accumulates for local Euclidean distance ($160\text{ operations}$)
  \item A 3-way minimum comparison and an addition ($4\text{ operations}$)
  \item Total arithmetic load: $250{,}000 \times 164 \approx 41 \times 10^6\text{ operations (41 MOPs)}$
\end{itemize}
On an embedded CPU, executing 250,000 sequential cell updates with nested memory lookups and branch checks consumes over $15\text{ ms}$---exceeding the real-time $10\text{ ms}$ streaming frame budget. Because all cells along any anti-diagonal wavefront ($i + j = k$) possess zero mutual data dependencies, this $250{,}000$-cell computational workload maps directly into a spatial systolic array accelerator on FPGA fabric, which we design and pipeline in Chapters~8 and~9.

\begin{groundbreaking}
Speech evaluation cannot be solved by a single algorithm. Statistical ASR is too forgiving, silently hallucinating correctness through language model priors; pitch correlation is too fragile, collapsing under non-uniform vowel stretching and falling blind to humming exploits. Robust edge evaluation demands a dual-tier multi-modal synthesis:
\begin{enumerate}\setlength{\itemsep}{0pt}
  \item \textbf{Tier 1 (Mel-DTW Elastic Alignment):} An 80-channel Log-Mel distance grid aligns physical vocal tract formants ($F_1, F_2, F_3$), instantly rejecting humming tricks and extracting the non-linear warping path $\mathcal{P}$.
  \item \textbf{Tier 2 (Warp-Guided Prosodic Correlation):} Pearson intonation ($F_0$) is evaluated strictly along $\mathcal{P}$, eliminating phase inversion and transforming a false negative failure ($r < 0$) into verified in-phase correlation ($r_{\text{warped}} > 0$).
\end{enumerate}
\end{groundbreaking}

\clearpage
% ==================== PAGE 7 ====================
\diagsec{Silicon and Prosodic Damage Board}

\noindent\begin{minipage}{\textwidth}
\centering
\resizebox{0.95\textwidth}{!}{\input{chapter01/fig_1_6a.tex}}\\[4pt]
{\scriptsize\textbf{Figure 1.6a}: \textbf{The pronunciation evaluation alignment damage board.} (1)~\textbf{What it is:} Comparison of temporal alignment paradigms between a hesitant learner (``Gooood... morning'', $1.8\text{ s}$) and an instructor reference (``Good morning'', $1.0\text{ s}$). (2)~\textbf{What to look at:} In Panel~(a), uniform linear resampling compresses the prolonged vowel across the word boundary, forcing frame 45 to multiply rising intonation against falling intonation ($\Delta x \cdot \Delta y < 0$), causing catastrophic phase inversion ($r < 0$). In Panel~(b), Dynamic Time Warping (DTW) on 80-d Log-Mel features elastically absorbs vowel stretching via non-linear path $\mathcal{P}$, preserving in-phase pitch correlation ($r_{\text{warped}} > 0$). (3)~\textbf{What it proves:} Accurate speech evaluation demands multi-modal decoupling: spectral features lock physical phonetic boundaries, while pitch correlation measures intonation strictly along the elastic path.\par}
\end{minipage}

\vspace{8pt}
\textbf{Organic Bridge to Chapter 2: The Transition from Algorithmic Formulation to Hardware Description.} Across Chapter~1, we deconstructed the mathematics, acoustic physics, and silicon failure modes of edge voice processing. We grounded the STFT and Mel filterbank in physical Block RAM primitives (Section~1.4), proved why standard parameters form an immovable Pareto boundary (Section~1.5), and established why speech evaluation requires dual-tier elastic dynamic programming (Section~1.6). 

However, mathematical equations and algorithmic flowcharts cannot be directly synthesized onto physical FPGA silicon. In Chapter~2, we transition from high-level algorithms to hardware description languages (Verilog/SystemVerilog and VHDL). We explore how continuous time and floating-point real numbers are discretized into synchronous clock domains, fixed-point integer representations, and pipelined register transfers, constructing the foundational digital logic blocks that breathe physical life into our streaming voice architecture.

\end{document}